% GENERATED FILE. Edit canonical lessons, then run npm run generate:books.
% canonical-source-hash: fnv1a64:fa52222c17709045
% canonical-lessons: JA-C140-senshuu, JA-C140-sengetsu, JA-C140-kyonen, JA-C140-sakki, JA-C140-owarimashita

\chapter{Last Week, Last Month, Last Year: the Past}
\label{ch:ja-last-week-last-month-last-year-the-past}
\hlchaptermodality{140}

\begin{chapteropening}
\textbf{By the end of this chapter:} \emph{I can say what I did and did not do, and when: last week, last month, last year, a moment ago.}

Say what you did and did not do with -mashita and -masen deshita, placing each with senshuu, sengetsu, kyonen or sakki, and end with owarimashita.
\end{chapteropening}

\section[\texorpdfstring{\ja{せんしゅう} (senshū)}{せんしゅう (senshū)}]{\ja{せんしゅう} (senshū) — last week}

\label{lesson:JA-C140-senshuu}



\begin{quote}
Four recalls before the new word.

\begin{itemize}
\raggedright
  \item \emph{Recall:} say \emph{nowhere}, with a negative verb — \textbf{R1}, one lesson back
  \item \emph{Recall:} answer \emph{no, that is not right} — \textbf{R2}, five lessons back
  \item \emph{Recall:} say \emph{a pencil} — \textbf{R3}, twenty lessons back
  \item \emph{Recall:} ask for \emph{tea, please}, with the particle \textbf{\ja{を}} — \textbf{R4}, eighty lessons back
\end{itemize}
\end{quote}

\subsection*{You'll want to know: \ja{せんしゅう}}
\textbf{\ja{せんしゅう}} — \emph{senshū} — \textquotedblleft{}last week\textquotedblright{}. It sits beside \textbf{\ja{こんしゅう}}, \emph{konshū}, \textquotedblleft{}this week\textquotedblright{}, which you already know.

To say what happened, the polite verb changes \textbf{\ja{ます}} to \textbf{\ja{ました}}, \emph{mashita}:

\begin{quote}
\textbf{\ja{せんしゅう} \ja{はたらきました。}} — \emph{senshū hatarakimashita} — \textquotedblleft{}I worked last week.\textquotedblright{}
\end{quote}

\textbf{\ja{はたらきます}} $\to$ \textbf{\ja{はたらきました}}, \textbf{\ja{ねます}} $\to$ \textbf{\ja{ねました}}. Every polite verb makes its past this way.

\subsection*{Your turn}
\begin{itemize}
\raggedright
  \item \emph{Say it:} \emph{senshū}
  \item \emph{Say it:} \emph{senshū hatarakimashita} — I worked last week
  \item \emph{Recall:} say \emph{I go home}, then \emph{I went home}: \emph{kaerimasu}, \emph{kaerimashita}
\end{itemize}

\begin{tcolorbox}[breakable,colback=teal!4,colframe=teal!35!black,title={Before you move on}]
What is \textquotedblleft{}last week\textquotedblright{}? (\textbf{\ja{せんしゅう}}, \emph{senshū}.) Say it once more.
\end{tcolorbox}
\section[\texorpdfstring{\ja{せんげつ} (sengetsu)}{せんげつ (sengetsu)}]{\ja{せんげつ} (sengetsu) — last month}

\label{lesson:JA-C140-sengetsu}



\begin{quote}
Four recalls before the new word.

\begin{itemize}
\raggedright
  \item \emph{Recall:} say \emph{last week} — \textbf{R1}, one lesson back
  \item \emph{Recall:} say \emph{is not}, plainly — \textbf{R2}, five lessons back
  \item \emph{Write it:} \textbf{\ja{ぴ}} from memory — \textbf{R3}, twenty lessons back
  \item \emph{Write it:} \textbf{\ja{を}} from memory — \textbf{R4}, eighty lessons back
\end{itemize}
\end{quote}

\subsection*{You'll want to know: \ja{せんげつ}}
\textbf{\ja{せんげつ}} — \emph{sengetsu} — \textquotedblleft{}last month\textquotedblright{}. The \textbf{\ja{せん}} at its start is the same as in \textbf{\ja{せんしゅう}}: \textquotedblleft{}the one before\textquotedblright{}.

\begin{quote}
\textbf{\ja{せんげつ} \ja{うみに} \ja{いきました。}} — \emph{sengetsu umi ni ikimashita} — \textquotedblleft{}I went to the sea last month.\textquotedblright{}
\end{quote}

\textbf{\ja{いく}}, \emph{iku}, \textquotedblleft{}to go\textquotedblright{}, gives \textbf{\ja{いきます}} and then \textbf{\ja{いきました}}.

\subsection*{Your turn}
\begin{itemize}
\raggedright
  \item \emph{Say it:} \emph{sengetsu}
  \item \emph{Say it:} \emph{sengetsu umi ni ikimashita} — I went to the sea last month
  \item \emph{Recall:} say \emph{last week}, then \emph{last month}
\end{itemize}

\begin{tcolorbox}[breakable,colback=teal!4,colframe=teal!35!black,title={Before you move on}]
What is \textquotedblleft{}last month\textquotedblright{}? (\textbf{\ja{せんげつ}}, \emph{sengetsu}.) Say it once more.
\end{tcolorbox}
\section[\texorpdfstring{\ja{きょねん} (kyonen)}{きょねん (kyonen)}]{\ja{きょねん} (kyonen) — last year}

\label{lesson:JA-C140-kyonen}



\begin{quote}
Three recalls before the new word.

\begin{itemize}
\raggedright
  \item \emph{Recall:} say \emph{last month} — \textbf{R1}, one lesson back
  \item \emph{Recall:} say \emph{nothing}, with a negative verb — \textbf{R2}, five lessons back
  \item \emph{Recall:} say \emph{one small animal} — \textbf{R3}, twenty lessons back
\end{itemize}
\end{quote}

\subsection*{You'll want to know: \ja{きょねん}}
\textbf{\ja{きょねん}} — \emph{kyonen} — \textquotedblleft{}last year\textquotedblright{}. It sits beside \textbf{\ja{ことし}}, \emph{kotoshi}, \textquotedblleft{}this year\textquotedblright{}.

To say what did NOT happen, the polite negative \textbf{\ja{ません}} adds \textbf{\ja{でした}}, \emph{deshita}: \textbf{\ja{ませんでした}}, \emph{masen deshita}.

\begin{quote}
\textbf{\ja{きょねん} \ja{どこにも} \ja{いきませんでした。}} — \emph{kyonen doko ni mo ikimasen deshita} — \textquotedblleft{}Last year I didn't go anywhere.\textquotedblright{}
\end{quote}

\noindent
\begin{tabularx}{\linewidth}{@{}>{\raggedright\arraybackslash}X>{\raggedright\arraybackslash}X>{\raggedright\arraybackslash}X@{}}
\toprule
\textbf{form} & \textbf{now, or as a plan} & \textbf{before now} \\
\midrule
affirmative & \textbf{\ja{いきます}} & \textbf{\ja{いきました}} \\
negative & \textbf{\ja{いきません}} & \textbf{\ja{いきませんでした}} \\
\bottomrule
\end{tabularx}

\subsection*{Your turn}
\begin{itemize}
\raggedright
  \item \emph{Say it:} \emph{kyonen}
  \item \emph{Say it:} \emph{kyonen doko ni mo ikimasen deshita} — Last year I didn't go anywhere
  \item \emph{Recall:} say \emph{I didn't work}: \emph{hatarakimasen deshita}
\end{itemize}

\begin{tcolorbox}[breakable,colback=teal!4,colframe=teal!35!black,title={Before you move on}]
What is \textquotedblleft{}last year\textquotedblright{}? (\textbf{\ja{きょねん}}, \emph{kyonen}.) Say it once more.
\end{tcolorbox}
\section[\texorpdfstring{\ja{さっき} (sakki)}{さっき (sakki)}]{\ja{さっき} (sakki) — a moment ago, a little while ago}

\label{lesson:JA-C140-sakki}



\begin{quote}
Four recalls before the new word.

\begin{itemize}
\raggedright
  \item \emph{Recall:} say \emph{last year} — \textbf{R1}, one lesson back
  \item \emph{Recall:} say \emph{nobody}, with a negative verb — \textbf{R2}, five lessons back
  \item \emph{Recall:} say \emph{a ticket} — \textbf{R3}, twenty lessons back
  \item \emph{Recall:} say \emph{there}, near the listener — \textbf{R4}, eighty lessons back
\end{itemize}
\end{quote}

\subsection*{You'll want to know: \ja{さっき}}
\textbf{\ja{さっき}} — \emph{sakki} — \textquotedblleft{}a moment ago\textquotedblright{}, \textquotedblleft{}a little while ago\textquotedblright{}: earlier today, not long before now. The small \textbf{\ja{っ}} holds the beat before \emph{ki}.

\begin{quote}
\textbf{\ja{さっき} \ja{おきました。}} — \emph{sakki okimashita} — \textquotedblleft{}I got up a little while ago.\textquotedblright{}
\end{quote}

\subsection*{Your turn}
\begin{itemize}
\raggedright
  \item \emph{Say it:} \emph{sakki}
  \item \emph{Say it:} \emph{sakki okimashita} — I got up a little while ago
  \item \emph{Recall:} say \emph{last year}, then \emph{a moment ago}, and say which is closer to now
\end{itemize}

\begin{tcolorbox}[breakable,colback=teal!4,colframe=teal!35!black,title={Before you move on}]
What is \textquotedblleft{}a moment ago\textquotedblright{}? (\textbf{\ja{さっき}}, \emph{sakki}.) Say it once more.
\end{tcolorbox}
\section[\texorpdfstring{\ja{おわりました} (owarimashita)}{おわりました (owarimashita)}]{\ja{おわりました} (owarimashita) — it finished, it is over (polite past)}

\label{lesson:JA-C140-owarimashita}



\begin{quote}
Four recalls before the new word.

\begin{itemize}
\raggedright
  \item \emph{Recall:} say \emph{a moment ago} — \textbf{R1}, one lesson back
  \item \emph{Recall:} say \emph{nowhere}, with a negative verb — \textbf{R2}, five lessons back
  \item \emph{Write it:} \textbf{\ja{ぷ}} from memory — \textbf{R3}, twenty lessons back
  \item \emph{Write it:} \textbf{\ja{そ}} from memory — \textbf{R4}, eighty lessons back
\end{itemize}
\end{quote}

\subsection*{You'll want to know: \ja{おわりました}}
\textbf{\ja{おわりました}} — \emph{owarimashita} — \textquotedblleft{}it finished\textquotedblright{}, \textquotedblleft{}it is over\textquotedblright{}. It is the polite past of \textbf{\ja{おわる}}, \emph{owaru}, \textquotedblleft{}to come to an end\textquotedblright{}, which follows the \emph{-u} pattern: \textbf{\ja{る}} becomes \textbf{\ja{り}}, then \textbf{\ja{ました}}.

\begin{quote}
\textbf{\ja{しごとが} \ja{おわりました。}} — \emph{shigoto ga owarimashita} — \textquotedblleft{}Work is over; I have finished work.\textquotedblright{}
\end{quote}

Now you can say what you did, and when:

\begin{quote}
\textbf{\ja{せんしゅう} \ja{はたらきました。}} — \textquotedblleft{}I worked last week.\textquotedblright{}
\end{quote}

\begin{quote}
\textbf{\ja{きょねん} \ja{どこにも} \ja{いきませんでした。}} — \textquotedblleft{}Last year I didn't go anywhere.\textquotedblright{}
\end{quote}

\subsection*{Your turn}
\begin{itemize}
\raggedright
  \item \emph{Say it:} \emph{owarimashita}
  \item \emph{Say it:} \emph{shigoto ga owarimashita} — Work is over
  \item \emph{Recall:} say what you did: \emph{senshū} (last week), \emph{sengetsu} (last month), \emph{kyonen} (last year) and \emph{sakki} (a moment ago), each with a verb in \emph{-mashita} or \emph{-masen deshita}, and end with \emph{owarimashita}
\end{itemize}

\begin{tcolorbox}[breakable,colback=teal!4,colframe=teal!35!black,title={Before you move on}]
Last week? (\textbf{\ja{せんしゅう}}, \emph{senshū}.) Last month? (\textbf{\ja{せんげつ}}, \emph{sengetsu}.) Last year? (\textbf{\ja{きょねん}}, \emph{kyonen}.) A moment ago? (\textbf{\ja{さっき}}, \emph{sakki}.) It is over? (\textbf{\ja{おわりました}}, \emph{owarimashita}.)
\end{tcolorbox}
